% chapters/ch5_study3.tex — 第 5 章 研究内容三
%
% 研究内容三：动态场景高斯地图的时间感知致密化（实证表征）。
% 章标题沿 OUTLINE 第四章结构定名；内容按 2026-11-07 用户拍板的「戊·两步走」
% 重定向为实证表征（monogs-ours paper3/ 一页纸 onepager_wu_boundary_2026-11-07 §一）。
%
% 数字来源纪律（WORKFLOW 铁律 2 数字零手抄 + WRITING_RULES 六）：
% 全部实验数字出自 monogs-ours/paper3/results/evidence/ 的判决文档与 judge JSON，
% 各节以 % 注释给出来源指针；剂量过供给第三 seed 补批（Amendment A1）与 async
% 迁移检查（P3-D2-043）收口后相关数字随之刷新（2026-11-08 派发中）。
% 测量口径红线：仪器口径（单线程+确定性栅格化）读数与 async 部署口径读数
% 永不混引，绝对 ATE 不跨口径、不跨机并列。

\chapter{动态场景高斯地图的时间感知致密化}
\label{ch:study3}

\section{引言}

第 3 章处理的是观测进入系统的入口，第 4 章处理的是系统跑完之后留下的地图。本章回到
两者之间那个持续发生的环节：地图的生长。在线高斯 SLAM 的地图不是一次建成的，系统每
隔一段迭代就执行一次致密化检查，把梯度大的点一分为二或克隆一份，让地图跟着观测逐步
长出几何支撑。这套机制继承自静态场景的三维高斯泼溅，在动态场景里它面对的问题是同一个吗，
默认参数还需要外部干预吗，干预的收益由什么决定，这三个问题构成本章的研究对象。
% 来源：paper3/PROBLEM.md §1-2（研究问题）；章节边界句沿 ch3/ch4 入口-地图-生长的组织线
% （OUTLINE.md §二 Q2：可靠观测 → 地图表示与维护 → 地图生长）。

对这三个问题，本章的回答是一份实证表征。选择轴上，我们系统测试了五类外置干预机制，
包括按统计纯度加权、按观测次数设资格门、按未命中计数加剪枝、以及两种剂量调整，
没有一类产生净收益，管线自带的剪枝自调节承担了原本指望干预承担的职能。剂量轴上，
测量给出从完全关闭到四倍阈值的三测点前沿，完全关闭使跟踪误差恶化三点六厘米、
地图缩到原来的百分之十五，说明致密化是明确的正资产，同时默认剂量也不在前沿的
最低点。序列之间，同样的减量干预在人群序列回收零点七厘米，在移动箱体序列反而
有害，收益方向不由场景的动态占比预测。这些读数合起来给出一个可操作的判据，
告诉后来者在什么条件下不必再设计此类机制，也把算力指向真正有杠杆的地方。
% 来源：onepager_wu_boundary_2026-11-07 §一（主张，2026-11-07 按泛化探针收窄）；
% 判决链 dose_a4_instrument / dose_scan / dose_oversupply（2026-11-06~08）。
% 注：正文中数字为骨架期锚定值，第三 seed 补批收口后刷新（Amendment A1，2026-11-08）。

\section{致密化机理与默认行为}
\label{sec:study3-mechanism}

理解干预为什么无效，先要看清管线在没有干预时做了什么。系统每积累一百次迭代就执行
一轮致密化检查：先按位置梯度挑出一批候选点，符合条件的高斯被克隆或分裂，使地图在
欠支撑的区域加出新的基元；随后按不透明度做一轮剪枝，把长期得不到观测支撑的基元
直接删除。clone 与 split 的触发由一个梯度阈值控制，阈值越低、每轮长出的基元越多，
这就是本文所说的剂量。剪枝不设阈值，凡不透明度衰减到下限以下的基元一律删除，
它是管线自带的清除机制。
% 来源：slam_backend.py densify_and_prune 调用点与 combined_prune/method.yaml
% 默认参数（densify_grad_threshold 0.0002、percent_dense 0.01、每 100 iter 执行）；
% 机理描述与 paper3/PROBLEM.md §1 一致。

这套默认行为里最关键的事实是剪枝自调节的强度。对一个生长中的地图连续记账的结果是，
致密化累计新增的基元里只有不到两成会被剪枝删除，其余都存活下来成为地图的长期成员；
换句话说，管线并非无节制地生长，无支撑的基元在出生后一两个检查周期内就会消亡。
这解释了后续各节一个反复出现的现象：外置干预试图替管线决定在哪里长、长多少，
而管线自身的选择与清除机制已经承担了大部分职能，干预要产生净收益就必须先超过
这个本底。
% 来源：paper3/results/evidence/growth_waste_balance_2026-11-06.md（prune÷新增 =
% 16.7%，83% 新生基元存活，40 帧仪器口径账本）；PROBLEM.md H1/H2 在本管线未获支持的
% 四条独立证据链之第一条（P3-D2-036）。

\section{测量方法}
\label{sec:study3-instrument}

本章的结论大多以零点几厘米的效应量出现，而跟踪误差本身就有厘米级的波动，测量装置
因此成为第一个要解决的问题。默认部署配置里有两个与设备无关的不确定性来源。其一是
异步调度，前端与后端各自推进，后端在两次关键帧之间跑多少次优化由墙钟决定，同样的
配置两次运行会走出不同的迭代数；其二是栅格化反向传播里的跨像素原子加法，浮点加法
的顺序不确定，梯度带着微小差异进入后续每一步。两者叠加的实测后果是，同一配置同一
随机种子重跑一次，跟踪误差可以相差两厘米以上，地图规模可以差出近一倍。
% 来源：device_floor_3090_2026-11-05.md（同 config 同 seed 重跑 online 极差 +96.8%，
% ATE 自差最高 2.14cm）；detraster_residual_bisect_2026-11-05.md（async 自由跑为
% 设备无关分叉源，同 20 帧两跑 735 vs 700 个不同 iteration）；P3-D2-022~026 装置诊断链。

测量方法把这两个来源逐个关闭。栅格化反向传播的浮点原子加法改写为整数镜像累加，
加法顺序不再影响结果；前端与后端都切到单线程固定步长，迭代数不再依赖墙钟；在此
之上叠加数值确定性套件。关闭之后，同一配置同一随机种子的两次运行在全部判据文件上
逐位一致，这一性质称为仪器有效性，作为每个正式批的前置门，两次运行有任何一位
不同则整批作废。这把装置的自差从厘米级压到零，使零点五厘米级的机制效应第一次
可以与噪声区分。
% 来源：prereg_instrument_mode_2026-11-05.md（三项锁定开关 + 三环境 + R-valid 判据）；
% detraster_L2_2x2_2026-11-05.md（2×2 四格，仅「单线程 × det-ON」达逐位一致）。

仪器口径与部署口径是两把不同的尺子，读数不能混用。单线程关闭了后端持续优化带来的
地图新鲜度，两者的跟踪误差与地图规模系统性地不同，本章所有臂级比较都在仪器口径内
完成；部署口径下的对应表现单列检验，两套读数分别登记，绝对值不与文献表格并列。
测量方法本身是使本章结论可判的前提，属于配套与评测卫生，不作为贡献主张。
% 来源：prereg_instrument_mode §0「它不是部署配置」（ATE 1.168 vs async 1.177、
% 高斯 19771 vs 25976）；DISCIPLINE「测量口径」#3 口径单源化 + 论文叙事纪律
%（评测协议不进贡献栏，用户 2026-09-25 裁定）。


